%% Application note template — Singlet AI
%% Short-form (2--4 pages) describing a software tool or database.
\documentclass[10pt,twocolumn,letterpaper]{article}
\usepackage{singletai-preprint}
\usepackage{graphicx}
\usepackage{amsmath}
\usepackage{booktabs}
\usepackage{listings}
\usepackage{natbib}
\bibliographystyle{unsrtnat}

% Code listing style
\lstset{
  basicstyle=\ttfamily\footnotesize,
  columns=fullflexible,
  breaklines=true,
  frame=single,
  backgroundcolor=\color[gray]{0.96},
  language=Python,
  showstringspaces=false,
}

% ── Metadata ─────────────────────────────────────────────────────
\title{Software Name: A Brief Description}
\author{
  Zach DeBruine\textsuperscript{1,*}
  \\[6pt]
  \textsuperscript{1}Grand Valley State University, Allendale, MI
}

\begin{document}
\maketitle
\correspondingauthor{debruinz@gvsu.edu}

\begin{abstract}
  % 50--100 words. What does the tool do, why does it matter?
  Abstract text goes here.
\end{abstract}

\keywords{software \and bioinformatics \and single-cell}

\section{Introduction}
% Brief motivation (1--2 paragraphs).

\section{Features}
% Key capabilities, with code examples:
% \begin{lstlisting}
% import singlet
% adata = singlet.load("GSE136831")
% \end{lstlisting}

\section{Implementation}
% Architecture, dependencies, performance.

\section{Availability}
Source code: \url{https://github.com/Singlet-AI/}.
Documentation: \url{https://singlet-ai.github.io/}.

\bibliography{refs}

\end{document}
